\documentclass[a4paper,13pt,3p,twoside]{report}
\usepackage{scrextend}
\changefontsizes{13pt}
\usepackage[utf8]{vietnam}
\usepackage[top=2cm, bottom=2cm, left=3.5cm, right=2.5cm]{geometry}
\usepackage{xurl}
\usepackage{appendix}
\usepackage{babel}
\usepackage{xcolor}
\usepackage{tikz}
\usetikzlibrary{positioning,arrows.meta}
\usepackage{outlines}
\usepackage{graphicx} % Cho phép chèn hỉnh ảnh
\usepackage{fancybox} % Tạo khung box
\usepackage{indentfirst} % Thụt đầu dòng ở dòng đầu tiên trong đoạn
\usepackage{amsthm} % Cho phép thêm các môi trường định nghĩa
\usepackage{latexsym} % Các kí hiệu toán học
\usepackage{amsmath} % Hỗ trợ một số biểu thức toán học
\usepackage{amssymb} % Bổ sung thêm kí hiệu về toán học
\usepackage{amsbsy} % Hỗ trợ các kí hiệu in đậm
\usepackage{times} % Chọn font Time New Romans
\usepackage{array} % Tạo bảng array

\usepackage{enumitem} % Cho phép thay đổi kí hiệu của list
\usepackage{subfiles} % Chèn các file nhỏ, giúp chia các chapter ra nhiều file hơn
\usepackage{titlesec} % Giúp chỉnh sửa các tiêu đề, đề mục như chương, phần,..
\usepackage{titletoc}
\usepackage{chngcntr} % Dùng để thiết lập lại cách đánh số caption,..
\usepackage{pdflscape} % Đưa các bảng có kích thước đặt theo chiều ngang giấy
\usepackage{afterpage}
\usepackage{algorithm} % Cung cấp môi trường float cho algpseudocode
\usepackage{algpseudocode}
\floatname{algorithm}{Thuật toán}
\algrenewcommand\algorithmicrequire{\textbf{Input:}}
\algrenewcommand\algorithmicensure{\textbf{Output:}}
\algrenewcommand\algorithmiccomment[1]{\hfill\(\triangleright\) #1}
\usepackage{capt-of} % Cho phép sử dụng caption lớn đối với landscape page
\usepackage{multirow} % Merge cells
\usepackage{fancyhdr} % Cho phép tùy biến header và footer

% \usepackage[natbib,backend=biber,style=ieee]{biblatex} % Giúp chèn tài liệu tham khảo

\usepackage[font=small,labelfont=bf]{caption}

\usepackage{listings}
\usepackage{float}
\usepackage{subcaption}

\usepackage[nonumberlist, nopostdot, nogroupskip, acronym]{glossaries}
\usepackage{glossary-superragged}
\setglossarystyle{superraggedheaderborder}
\usepackage{setspace}
\usepackage{parskip}

% package content table
\usepackage{tocbasic}

\usepackage{blindtext}
\usepackage{booktabs}
\usepackage{makecell}
\usepackage{hyperref}
\usepackage{makecell}
\usepackage{tabularx}

% ===================================================

\renewcommand{\bibname}{Tai_lieu_tham_khao} 
\usepackage[backend=biber,style=ieee]{biblatex}  %backend=biber is 'better'

\addbibresource{Tai_lieu_tham_khao.bib} % chèn file chứa danh mục tài liệu tham khảo vào 
\usepackage{xurl}                       % cho phép ngắt URL/DOI ở mọi ký tự
\setcounter{biburlnumpenalty}{100}
\setcounter{biburlucpenalty}{100}
\setcounter{biburllcpenalty}{100}
\appto{\bibsetup}{\sloppy}              % nới lỏng căn lề cho phần tài liệu tham khảo
\newcolumntype{Y}[1]{>{\RaggedRight\arraybackslash}p{#1\dimexpr\textwidth-6\tabcolsep\relax}}
\newcolumntype{C}[1]{>{\Centering\arraybackslash}p{#1\dimexpr\textwidth-6\tabcolsep\relax}}
\newcolumntype{L}[1]{>{\raggedright\arraybackslash}p{#1}}

\usepackage{listing} % Phần này cho phép chèn code và formatting code như C, C++, Python
\lstset{
    basicstyle=\ttfamily\small,
    breaklines=true,              % tự xuống dòng khi quá dài
    breakatwhitespace=false,      % cho phép ngắt cả giữa từ/ký tự (cần cho tên dài, chuỗi dài)
    columns=fullflexible,         % tránh khoảng trắng bị giãn khi ngắt dòng
    keepspaces=true,
    prebreak=\mbox{\textcolor{gray}{\tiny$\searrow$}},     % dấu hiệu ở cuối dòng bị ngắt
    postbreak=\mbox{\textcolor{gray}{\tiny$\hookrightarrow$}\space}, % dấu hiệu đầu dòng tiếp theo
    frame=single,
    xleftmargin=1.5em,            % chừa chỗ cho số dòng
    framexleftmargin=1.5em,
    numbers=left,
    numberstyle=\tiny\color{gray},
    showstringspaces=false,
    tabsize=4
}

\definecolor{backcolour}{RGB}{248,248,245}
\definecolor{codegray}{RGB}{128,128,128}
\definecolor{codegreen}{RGB}{0,128,0}
\definecolor{codepurple}{RGB}{136,0,136}
\definecolor{codeblue}{RGB}{0,0,200}
\definecolor{codeorange}{RGB}{200,90,0}

\lstdefinestyle{treestyle}{
    backgroundcolor=\color{backcolour},
    basicstyle=\ttfamily\scriptsize\color{black},
    breaklines=true,
    breakatwhitespace=false,
    columns=fullflexible,
    numbers=none,
    frame=single,
    framerule=0.5pt,
    rulecolor=\color{codegray},
    keepspaces=true,
    tabsize=2,
    commentstyle=\color{codegreen}\itshape,
    morecomment=[l]{\#},
    literate=
      {├──}{{|-- }}4
      {└──}{{\textasciigrave-- }}4
      {│}{{|}}1
      {─}{{-}}1
}

\lstdefinestyle{pystyle}{
    language=Python,
    backgroundcolor=\color{backcolour},
    basicstyle=\ttfamily\small\color{black},
    keywordstyle=\color{codeblue}\bfseries,       % def, for, in, return...
    keywordstyle=[2]\color{codepurple},           % hàm built-in (nếu khai báo)
    commentstyle=\color{codegreen}\itshape,
    stringstyle=\color{codeorange},
    numberstyle=\tiny\color{codegray},
    numbers=left,
    numbersep=6pt,
    frame=single,
    rulecolor=\color{codegray},
    breaklines=true,
    breakatwhitespace=false,
    columns=fullflexible,
    keepspaces=true,
    showstringspaces=false,
    tabsize=4,
    xleftmargin=1.5em,
    framexleftmargin=1.5em,
    morekeywords={self, True, False, None},
    morekeywords=[2]{len, range, list, bytes, int, ntt, intt, encode, decode, compress, decompress}
}

\input{Tu_viet_tat}

% ===================================================


\fancypagestyle{plain}{%
\fancyhf{} % clear all header and footer fields
\fancyfoot[RO,RE]{\thepage} %RO=right odd, RE=right even
\renewcommand{\headrulewidth}{0pt}
\renewcommand{\footrulewidth}{0pt}}

\setlength{\headheight}{10pt}

\def \TITLE{ĐỒ ÁN TỐT NGHIỆP}
\def \AUTHOR{Trần Văn A}

% ===================================================
\titleformat{\chapter}[hang]{\centering\bfseries}{CHƯƠNG \thechapter.\ }{0pt}{}[]

\titleformat 
    {\chapter} % command
    [hang] % shape
    {\centering\bfseries} % format
    {CHƯƠNG \thechapter.\ } % label
    {0pt} %sep
    {} % before
    [] % after
\titlespacing*{\chapter}{0pt}{-20pt}{20pt}

\titleformat
    {\section} % command
    [hang] % shape
    {\bfseries} % format
    {\thechapter.\arabic{section}\ \ \ \ } % label
    {0pt} %sep
    {} % before
    [] % after
\titlespacing{\section}{0pt}{\parskip}{0.5\parskip}

\titleformat
    {\subsection} % command
    [hang] % shape
    {\bfseries} % format
    {\thechapter.\arabic{section}.\arabic{subsection}\ \ \ \ } % label
    {0pt} %sep
    {} % before
    [] % after
\titlespacing{\subsection}{30pt}{\parskip}{0.5\parskip}

\renewcommand\thesubsubsection{\alph{subsubsection}}
\titleformat
    {\subsubsection} % command
    [hang] % shape
    {\bfseries} % format
    {\alph{subsubsection}, \ } % label
    {0pt} %sep
    {} % before
    [] % after
\titlespacing{\subsubsection}{50pt}{\parskip}{0.5\parskip}

% \newcommand{\titlesize}{\fontsize{18pt}{23pt}\selectfont}
% \newcommand{\subtitlesize}{\fontsize{16pt}{21pt}\selectfont}
% \titleclass{\part}{top}
% \titleformat{\part}[display]
%   {\normalfont\huge\bfseries}{\centering}{20pt}{\Huge\centering}
% \titlespacing{\part}{0pt}{em}{1em}
% \titlespacing{\section}{0pt}{\parskip}{0.5\parskip}
% \titlespacing{\subsection}{0pt}{\parskip}{0.5\parskip}
% \titlespacing{\subsubsection}{0pt}{\parskip}{0.5\parskip}



% ===================================================
\usepackage{hyperref}
\hypersetup{pdfborder = {0 0 0}}
\hypersetup{pdftitle={\TITLE},
	pdfauthor={\AUTHOR}}
	
\usepackage[all]{hypcap} % Cho phép tham chiếu chính xác đến hình ảnh và bảng biểu

\graphicspath{{figures/}{../figures/}} % Thư mục chứa các hình ảnh

\counterwithin{figure}{chapter} % Đánh số hình ảnh kèm theo chapter. Ví dụ: Hình 1.1, 1.2,..

\title{\bf \TITLE}
\author{\AUTHOR}

\setcounter{secnumdepth}{3} % Cho phép subsubsection trong report
% \setcounter{tocdepth}{3} % Chèn subsubsection vào bảng mục lục

\theoremstyle{definition}
\newtheorem{example}{Ví dụ}[chapter] % Định nghĩa môi trường ví dụ
\newtheorem{theorem}{Định lý}[chapter] % Định nghĩa môi trường ví dụ
\newtheorem{lemma}{Bổ đề}[chapter] % Định nghĩa môi trường ví dụ

\onehalfspacing
\setlength{\parskip}{6pt}
\setlength{\parindent}{15pt}

\theoremstyle{definition}
\newtheorem{dn}{Định nghĩa}[section]
\providecommand{\dnautorefname}{Định nghĩa}
\newtheorem{bt}{Bài toán}[chapter]
\providecommand*{\btautorefname}{Bài toán}
\newtheorem{vd}{\bfseries Ví dụ}[section]
\providecommand*{\mdautorefname}{Ví dụ}
\newtheorem{vd1}{\bfseries Ví dụ}[chapter]
\providecommand*{\mdautorefname}{Ví dụ}


% =========================== BODY ===============
\begin{document}
% \newgeometry{top=2cm, bottom=2cm, left=2cm, right=2cm}
\subfile{Bia} % Phần bìa
% \restoregeometry

% ===================================================
\pagestyle{empty} % Header và footer rỗng
%\newpage
%\pagenumbering{gobble} % Xóa page numbering ở cuối trang
%\subfile{chapters/0_1_subject.tex}

% \pagestyle{empty} % Header và footer rỗng
% \newpage
% \subfile{Chuong/0_0_Xac_nhan_chinh_sua.tex}
\newpage
\pagenumbering{gobble} % Xóa page numbering ở cuối trang
\subfile{Chuong/0_1_Loi_cam_on_Tom_tat.tex}

% % \pagestyle{empty} % Header và footer rỗng
% \newpage
% \pagenumbering{gobble} % Xóa page numbering ở cuối trang
% \subfile{Chuong/0_2_Mo_dau.tex}

% % \pagestyle{empty} % Header và footer rỗng
% \newpage
% \pagenumbering{gobble} % Xóa page numbering ở cuối trang
% \subfile{Chuong/0_4_Tom_tat_noi_dung_eng.tex}

% ===================================================
% \pagestyle{empty} % Header và footer rỗng
\newpage
\pagenumbering{gobble} % Xóa page numbering ở cuối trang
\renewcommand*\contentsname{MỤC LỤC}
\titlecontents{chapter}
    [0.0cm]             % left margin
    {\bfseries\vspace{0.3cm}}                  % above code
    {{\bfseries{\scshape} CHƯƠNG \thecontentslabel.\ }} % numbered format
    {}         % unnumbered format
    {\titlerule*[0.3pc]{.}\contentspage}         % filler-page-format, e.g dots
    
\titlecontents{section}
    [0.0cm]             % left margin
    {\vspace{0.3cm}}                  % above code
    {\thecontentslabel \ } % numbered format
    {}         % unnumbered format
    {\titlerule*[0.3pc]{.}\contentspage}         % filler-page-format, e.g dots
    
\titlecontents{subsection}
    [1.0cm]             % left margin
    {\vspace{0.3cm}}                  % above code
    {\thecontentslabel \ } % numbered format
    {}         % unnumbered format
    {\titlerule*[0.3pc]{.}\contentspage}         % filler-page-format, e.g dots

 % Tạo mục lục tự động
\addtocontents{toc}{\protect\thispagestyle{empty}}
\tableofcontents 
\thispagestyle{empty}
\cleardoublepage

\pagenumbering{roman}
%Tạo danh mục hình vẽ.
\renewcommand{\listfigurename}{DANH MỤC HÌNH VẼ}
{\let\oldnumberline\numberline
\renewcommand{\numberline}{Hình~\oldnumberline}
\listoffigures} 
% \phantomsection\addcontentsline{toc}{section}{\numberline {} DANH MỤC HÌNH VẼ}
\newpage


 %Tạo danh mục bảng biểu.
\renewcommand{\listtablename}{DANH MỤC BẢNG BIỂU}
{\let\oldnumberline\numberline
\renewcommand{\numberline}{Bảng~\oldnumberline}
\listoftables}
% \phantomsection\addcontentsline{toc}{section}{\numberline {} DANH MỤC BẢNG BIỂU}

\glsaddall 
% \renewcommand*{\glossaryname}{Danh sách thuật ngữ}
\renewcommand*{\acronymname}{DANH MỤC THUẬT NGỮ VÀ TỪ VIẾT TẮT}
\renewcommand*{\entryname}{Thuật ngữ}
\renewcommand*{\descriptionname}{Ý nghĩa}
\printnoidxglossaries
% \phantomsection\addcontentsline{toc}{section}{\numberline {} DANH MỤC THUẬT NGỮ VÀ TỪ VIẾT TẮT}

% \renewcommand\appendixname{PHỤ LỤC}
% \renewcommand\appendixpagename{PHỤ LỤC}
% \renewcommand\appendixtocname{PHỤ LỤC}
\widowpenalty=1000
\clubpenalty=1000
\raggedbottom


% ===================================================


\newpage
\pagenumbering{arabic}

\pagestyle{fancy}
\fancyhf{}
\fancyhead[RE, LO]{\leftmark}
%\fancyhead[LE]{\rightmark}
\fancyfoot[RE, LO]{\thepage}
% \pagestyle{empty} % Header và footer rỗng
\newpage
\subfile{Chuong/0_2_Mo_dau.tex}

% \chapter{GIỚI THIỆU ĐỀ TÀI}
% \subfile{Chuong/1_Gioi_thieu} % Phần mở đầu

\newpage
%\pagestyle{fancy} % Áp dụng header và footer
\chapter{CƠ SỞ LÝ THUYẾT}
\subfile{Chuong/1_Co_so_ly_thuyet}

\newpage
%\pagestyle{fancy} % Áp dụng header và footer
\chapter{PHƯƠNG PHÁP ĐỀ XUẤT}
\subfile{Chuong/2_Thiet_ke_mo_hinh}



\newpage
%\pagestyle{fancy} % Áp dụng header và footer
\chapter{TRIỂN KHAI VÀ ĐÁNH GIÁ AN TOÀN}
\label{chap:experiment}
\subfile{Chuong/3_Trien_khai_va_danh_gia}

% \newpage
% %\pagestyle{fancy} % Áp dụng header và footer
% \chapter{PHÂN TÍCH HIỆU NĂNG}
% \subfile{Chuong/5_Phan_tich_hieu_nang}

\newpage
\chapter*{KẾT LUẬN}
\subfile{Chuong/4_Ket_luan}

% \newpage
% \label{chapter:reference}
% \subfile{Chuong/7_Tai_lieu_tham_khao}

% ===================================================
\newpage
\renewcommand\bibname{TÀI LIỆU THAM KHẢO}
\printbibliography[heading=bibintoc]
% \phantomsection\addcontentsline{toc}{chapter}{TÀI LIỆU THAM KHẢO}

% \appendixpage
% \appendix
% \addappheadtotoc

\titleformat{\chapter}[hang]{\centering\bfseries}{ \thechapter.\ }{0pt}{}[]
\titlespacing*{\chapter}{0pt}{-20pt}{20pt}

\titlecontents{chapter}
    [0.0cm]             % left margin
    {\bfseries\vspace{0.3cm}}                  % above code
    {{\bfseries{\scshape} \thecontentslabel.\ }} % numbered format
    {}         % unnumbered format
    {\titlerule*[0.3pc]{.}\contentspage} 
    
% \chapter{HƯỚNG DẪN VIẾT ĐATN}
% \subfile{Chuong/Phu_luc_A}

% \newpage
% \chapter{ĐẶC TẢ USE CASE}
% \subfile{Chuong/Phu_luc_B}

\end{document}
